% ---
% id: meat-core/frank-wolfe/cor:fo-dropped
% kind: cor
% title: "First-order bound for the dropped program"
% summary: >-
%   In either source mode the first-order bound holds for the dropped program
%   ([[def:dropped]]), with maximal support asked only of key-active pairs
%   (acceptance value, joint sector) whose entropy is kept. It bounds h_low
%   from below, also at iterates feasible only up to windows.
% prominence: main
% status: implemented
% uses: [thm:firstorder, lem:fo-face, lem:fo-khyp, def:relaxed-feasible-set, prop:relaxed-set-properties, def:first-order-setting, meat-core/def:dropped, meat-core/prop:drop-bound, meat-core/cor:drop-program, meat-core/def:joint-sectors, meat-core/lem:joint-sectors, meat-core/lem:choi-factor, meat-core/thm:choi, meat-core/thm:gradient, meat-core/cor:full-gradient, meat-core/def:branch-support, meat-core/lem:branch-support, foundations/def:kl, meat-core/def:admissible-declaration, meat-core/def:admissible-acceptance-set, meat-core/def:feasible, meat-core/def:untrusted-choi]
% mentions: [meat-core/thm:convexity, meat-core/thm:branch-formula, meat-core/def:declaration]
% assumes: []
% refs: []
% provenance: newly-derived
% implements: certiqs_key.solvers.frank_wolfe.joint_driver:run_joint_frank_wolfe
% lean:
% notes: "The lower bounds of thm:firstorder (iv) are read here at the gradient of the dropped objective, a different linear form over the same sets; on the feasible side lem:feasible-side-dual supplies one for any linear form."
% ---

A run that gives up the deficits of a \tref{dropped set} takes its
first-order bound from the \tref{dropped program}. The corollary is to that
program what Theorem~\ref{thm:firstorder}\,(ii) and~(iiib) are to the
undropped one, with the support hypothesis asked of each kept key-active
joint pair and positivity asked of the dropped branch quantity, which is
never smaller than the undropped one. A \tref{fully dropped} joint sector
needs no maximal support, no rank decision and no matrix function. In the
relaxed reading the three hypotheses are conditions on the iterate alone
and name no set, and no quantity relaxed at $\varepsilon>0$ appears in the
conclusion: the chain
\eqref{eq:hlb-dropped} compares
$\underline{h}^{\mathrm{FO},\mathcal{D}}_{\alpha}$ against the exact
$\underline{h}^{\mathcal{D}}_{\alpha}$ directly, and no non-negativity of
a relaxed infimum is claimed. On an empty feasible set both values are
$+\infty$, so the bound needs the non-empty feasible set of an admissible
declaration only to be informative. The corollary holds in either source
mode: in \tref{untrusted-source mode} the Choi step is the identity, so the
frame and the operator coupling of the trusted-source proof become the
state blocks and the normalisation.

\begin{corollary}
\label{cor:fo-dropped}
Let $\alpha\in(1,\infty)$, let $\Decl$ satisfy Definition~\ref{def:admissible-declaration} in
the general format of Definition~\ref{def:declaration}, in either source
mode, in \tref{trusted-source mode} with a frame per Alice sector as in
Theorem~\ref{thm:choi}, let $\acSet$ satisfy
Definition~\ref{def:admissible-acceptance-set}, and let $\mathcal{D}$ be a dropped set
(Definition~\ref{def:dropped}). Write an Alice sector as
$\mathsf t\in\mathbb{T}$ in this claim, the letter $t$ indexing iterates as
in Theorem~\ref{thm:firstorder}. Put
$O^{J,\mathcal{D}}_{\alpha}(\mathbf q,J)\coloneqq
O^{\mathcal{D}}_{\alpha}(\mathbf q,\chi(J))$ and
$K^{J,\mathcal{D}}_{\hat c}\coloneqq K^{\mathcal{D}}_{\hat c}\circ\chi$, with
$\chi$ the assignment of that theorem in trusted-source mode. In
untrusted-source mode the program is read as
Definition~\ref{def:untrusted-choi} reads it, with no frame: $\chi$ is the
identity, $J=\omega$ is the list of the blocks
$J_{(\mathsf t,s)}=\omega^{(\mathsf t,s)}$ over the joint sectors,
$A_{0,\mathsf t}$ and $A_0$ stand for $A_{\mathsf t}$ and $A$, and
$\Feas^{J}$ is $\Feas$ of
Definition~\ref{def:feasible} in the general format, whose condition
coupling the blocks is the normalisation $\sum_{(\mathsf t,s)}
\tr\omega^{(\mathsf t,s)}=1$; so $O^{J,\mathcal{D}}_{\alpha}$ and
$K^{J,\mathcal{D}}_{\hat c}$ are $O^{\mathcal{D}}_{\alpha}$ and
$K^{\mathcal{D}}_{\hat c}$. Let
$x_t=(\mathbf q_t,J_t)\in\acSet\times\Feas^{J}$, $t\in\Itset$, satisfy
\begin{itemize}[leftmargin=3.4em]
\item[(H-POS$^{\mathcal{D}}$)] (H-POS) of Theorem~\ref{thm:firstorder} with
  $K^{J,\mathcal{D}}_{\hat c}$ in place of $K^{J}_{\hat c}$;
\item[(H-SUPP$^{\mathcal{D}}$)] for every $\hat c\in\Chkey$ and every joint
  pair $(\mathsf t,s)\in\mathbb{J}_K(\hat c)$ with
  $(\hat c,(\mathsf t,s))\notin\mathcal{D}$, the branch operator
  $\cG^{J}_{\hat c,s}(J_{t,(\mathsf t,s)})$ of the block
  $J_{t,(\mathsf t,s)}\in\Pos{A_{0,\mathsf t}\Qsys_s}$ attains the maximal
  support of that joint pair, the support it has at
  $J_{t,(\mathsf t,s)}\succ0$; nothing is asked of a pair that is key-inert
  or dropped;
\item[(H-NOPERT$^{\mathcal{D}}$)] $(\lmoq_t,\lmoJ_t)$ is the gradient of
  $O^{J,\mathcal{D}}_{\alpha}$ at $x_t$ itself, assembled as in (H-NOPERT)
  with the kept summands only and with $K^{J,\mathcal{D}}_{\hat c}$ in the
  scalar prefactors, its blocks on fully dropped joint sectors being
  \eqref{eq:drop-gradient} in Choi coordinates;
\end{itemize}
and let $(v_t,S_t)$ satisfy \eqref{eq:lmo} for these linear forms, or let,
at some or all $t$, lower bounds $b^{\mathrm A}_t$, $b^{\mathrm F}_t$ on
one or both of its minima, as in \eqref{eq:lmo-bounds} for these linear
forms, stand in for the pairings as Theorem~\ref{thm:firstorder}\,(iv)
lets them. Then, in
exact real arithmetic, with $\underline{h}^{\mathrm{FO},\mathcal{D}}_{\alpha}$
the expression \eqref{eq:hlb} for $O^{J,\mathcal{D}}_{\alpha}$, with those
replacements where they are made,
\begin{equation}\label{eq:hlb-dropped}
  \max\{0,\underline{h}^{\mathrm{FO},\mathcal{D}}_{\alpha}\}\;\le\;\underline{h}^{\mathcal{D}}_{\alpha}
  \;\le\;\hlow\;\le\;\hqkd(\CList).
\end{equation}
(H-POS$^{\mathcal{D}}$) is implied by (H-POS), and
(H-SUPP$^{\mathcal{D}}$) by (H-SUPP) read on joint blocks. For
$\mathcal{D}=\varnothing$ and $|\mathbb{T}|=1$ the statement is
Theorem~\ref{thm:firstorder}\,(ii), with its part~(iv) where lower bounds
stand in.

\emph{The relaxed reading.} Let $\varepsilon\ge0$ and let
$\Feas^{J}(\varepsilon)$ be the \tref{relaxed feasible set} of
Definition~\ref{def:relaxed-feasible-set}, read in the general format, with
one coupling clause per Alice sector in trusted-source mode and its
normalisation clause in untrusted-source mode. Let
$x_t\in\acSet\times\Feas^{J}(\varepsilon)$, $t\in\Itset$, satisfy
(H-POS$^{\mathcal{D}}$), (H-SUPP$^{\mathcal{D}}$) and
(H-NOPERT$^{\mathcal{D}}$) as stated, and let $(v_t,S_t)$ satisfy
\eqref{eq:lmo} \emph{as written}, or the lower bounds standing in for its
pairings satisfy \eqref{eq:lmo-bounds}, with the two minima over $\acSet$
and over the \emph{exact} $\Feas^{J}$. Then \eqref{eq:hlb-dropped} holds in
full, and the case $\varepsilon=0$ is the statement above.

\emph{An empty feasible set.} In place of the hypothesis on $\Decl$ above,
let $\Decl$ be a tuple of the form of Definition~\ref{def:declaration}, in
the general format and in either source mode, that satisfies conditions
(i)--(vi) and (viii) of Definition~\ref{def:admissible-declaration}, and form
$\Feas$, $O_{\alpha}$, $\hlow$ and $\underline{h}^{\mathcal{D}}_{\alpha}$ from
it by the formulas of Definitions~\ref{def:feasible} and~\ref{def:dropped}.
Let $\Feas=\varnothing$. Then
$\underline{h}^{\mathcal{D}}_{\alpha}=\hlow=+\infty$, so every real number is
a lower bound on both, the value of \eqref{eq:hlb} for
$O^{J,\mathcal{D}}_{\alpha}$ that a run forms at its iterates included.
Consequently the first two inequalities of \eqref{eq:hlb-dropped}, for a real
$\underline{h}^{\mathrm{FO},\mathcal{D}}_{\alpha}$, hold of every such
tuple, whether or not it satisfies condition~(vii): under the hypotheses above where it does,
and by the preceding sentence where it does not.
\end{corollary}

\begin{proof}
\emph{Idea.} The first inequality of \eqref{eq:hlb-dropped} follows the
four steps of the proof of Theorem~\ref{thm:firstorder}\,(ii) with the
changes below, preceded by the two things the dropped program needs that
that theorem's proof does not state: the support geometry of a joint pair
(Step~0) and the acceptance-support bookkeeping for $K^{J,\mathcal{D}}$
(Step~0$'$). The relaxed reading changes Step~2 only. The proof is
written once for both source modes, with the reading of untrusted-source
mode fixed in the paragraph on notation.

\emph{The last two inequalities} of \eqref{eq:hlb-dropped} are
Proposition~\ref{prop:drop-bound}\,(iii).

\emph{Notation.} Write $p=(\mathsf t,s)$ for a joint pair. In
trusted-source mode put
$W_{\mathsf t}\coloneqq(\tau^{(\mathsf t)}_A)^{1/2}\iota_{\mathsf t}$ and
$\chi_p(\cdot)\coloneqq(W_{\mathsf t}\otimes\idon{\Qsys_s})(\cdot)
(W_{\mathsf t}^{\dagger}\otimes\idon{\Qsys_s})$ for the assignment of
Theorem~\ref{thm:choi} in the general format, so that
$\chi=\bigoplus_p\chi_p$ and $W\coloneqq\bigoplus_{\mathsf t}W_{\mathsf t}$
satisfies $WW^{\dagger}=\tau_A$
(Lemma~\ref{lem:choi-factor}\,(ii)), and
$S_{\mathsf t}\coloneqq\operatorname{supp}\tau^{(\mathsf t)}_A
=\operatorname{ran}W_{\mathsf t}$
(Lemma~\ref{lem:choi-factor}\,(i)). In untrusted-source mode
(Definition~\ref{def:untrusted-choi}\,(a)) let $W_{\mathsf t}$ be the
inclusion of $\Hs{A_{\mathsf t}}$ into $\Hs{A}$ and $\chi_p$ the conjugation
by $W_{\mathsf t}\otimes\idon{\Qsys_s}$, that is, the inclusion of
$\Pos{A_{\mathsf t}\Qsys_s}$ into $\Pos{A\Qsys_s}$ as the subcone of
Definition~\ref{def:joint-sectors}; then $\chi=\bigoplus_p\chi_p$ is the
identity on lists over the joint sectors, and
$S_{\mathsf t}\coloneqq\Hs{A_{\mathsf t}}=\operatorname{ran}W_{\mathsf t}$.
In both modes $\chi_p$ is linear and preserves positive semidefiniteness, and
$\operatorname{ran}W_{\mathsf t}=S_{\mathsf t}\subseteq\Hs{A_{\mathsf t}}$.
Write
$\cG^{J}_{\hat c,p}\coloneqq\cG_{\hat c,p}\circ\chi_p$, which is the
restriction of $\cG^{J}_{\hat c,s}$ to the blocks supported on
$\Hs{A_{0,\mathsf t}\Qsys_s}$, the reading (H-SUPP$^{\mathcal{D}}$) writes,
and $X_{\hat c,p}\coloneqq\cG^{J}_{\hat c,p}(J_{t,p})$. Let
$\widetilde\Pi_{\mathsf t}$ be $\Pi^{A}_{\mathsf t}$ acting on the $A$ factor
of the branch space $\Hs{R\,XY\,C_AC_B\,A\Qsys_s}$, and put
\begin{equation}\label{eq:Wcs-joint}
  W_{\hat c,p}\coloneqq\widetilde\Pi_{\mathsf t}\,W_{\hat c,s}
  \qquad\text{with }W_{\hat c,s}\text{ of }\eqref{eq:Wcs}.
\end{equation}

\emph{Step 0 (the support geometry of a joint pair).} \emph{(a)} By
\eqref{eq:Wcs}, $W_{\hat c,s}=\sum_{\kappa\in\Kann_{\mathrm{pass}}(\hat c)}
K^{(s)}_{\kappa}(\Sadm\otimes\Hs{\Qsys_s})$, and
$\Sadm=\bigoplus_{\mathsf t}S_{\mathsf t}$. In trusted-source mode
$\Sadm=\operatorname{supp}\tau_A$, and
$\tau_A=\bigoplus_{\mathsf t}\tau^{(\mathsf t)}_A$ is block
diagonal along $\mathbb{T}$ by Definition~\ref{def:admissible-declaration}\,(iii) in the general
format, so $\Sadm=\bigoplus_{\mathsf t}\operatorname{supp}
\tau^{(\mathsf t)}_A$; in untrusted-source mode $\Sadm=\Hs{A}$
(Definition~\ref{def:branch-support}), which is
$\bigoplus_{\mathsf t}\Hs{A_{\mathsf t}}$ by the general format of
Definition~\ref{def:declaration}. By the computation in the proof of
Lemma~\ref{lem:joint-sectors}\,(i),
$K^{(s)}_{\kappa}(\Pi^{A}_{\mathsf t}\otimes\idon{})
=\widetilde\Pi_{\mathsf t}K^{(s)}_{\kappa}$, so each
$K^{(s)}_{\kappa}(S_{\mathsf t}\otimes
\Hs{\Qsys_s})$ lies in the range of $\widetilde\Pi_{\mathsf t}$; summing over
$\mathsf t$ and $\kappa$ exhibits $W_{\hat c,s}$ as a sum of subspaces of
\emph{mutually orthogonal} summands, hence
\begin{equation}\label{eq:Wcs-joint-split}
  W_{\hat c,s}=\bigoplus_{\mathsf t\in\mathbb{T}}W_{\hat c,(\mathsf t,s)},
  \qquad
  W_{\hat c,(\mathsf t,s)}=\sum_{\kappa\in\Kann_{\mathrm{pass}}(\hat c)}
  K^{(s)}_{\kappa}\big(S_{\mathsf t}\otimes
  \Hs{\Qsys_s}\big).
\end{equation}
\emph{(b)} For every Hermitian $\Delta$ on $\Hs{A_{0,\mathsf t}\Qsys_s}$,
$\operatorname{supp}\cG^{J}_{\hat c,p}(\Delta)\subseteq W_{\hat c,p}$. Indeed
$\Delta\oplus0$ is Hermitian on $\Hs{A_0\Qsys_s}$, so
Lemma~\ref{lem:fo-face}, \eqref{eq:fo-face} gives
$\operatorname{supp}\cG^{J}_{\hat c,s}(\Delta\oplus0)\subseteq W_{\hat c,s}$;
and $\chi_p(\Delta_{\pm})$ is supported on
$\Hs{A_{\mathsf t}\Qsys_s}$, $\operatorname{ran}W_{\mathsf t}$ lying in
$\Hs{A_{\mathsf t}}$, so by
Lemma~\ref{lem:joint-sectors}\,(i) the image of each lies in the
range of $\widetilde\Pi_{\mathsf t}$ and so does their difference. Intersecting
the two inclusions gives the claim by \eqref{eq:Wcs-joint}. \emph{(c)} For
$\varrho\succeq0$ the support of $\sum_{\kappa}K_{\kappa}\varrho
K_{\kappa}^{\dagger}$ is $\sum_{\kappa}K_{\kappa}(\operatorname{supp}\varrho)$
(the support of a sum, in the proof of Lemma~\ref{lem:branch-support}). At
$J_{t,p}\succ0$, $\operatorname{supp}\chi_p(J_{t,p})
=\operatorname{ran}(W_{\mathsf t}\otimes\idon{})
=S_{\mathsf t}\otimes\Hs{\Qsys_s}$, so
$\operatorname{supp}X_{\hat c,p}=W_{\hat c,p}$ by
\eqref{eq:Wcs-joint-split}, while for any $J_{t,p}\succeq0$ the inclusion
$\subseteq$ holds. So the ``maximal support of that joint pair'' of
(H-SUPP$^{\mathcal{D}}$) is $W_{\hat c,p}$, and that hypothesis reads
$\operatorname{supp}X_{\hat c,p}=W_{\hat c,p}$ for every kept key-active pair.

\emph{Step 0$'$ (the bookkeeping, which is Lemma~\ref{lem:fo-khyp} for
$K^{J,\mathcal{D}}$).} Part~(i) of that lemma is a statement about $\acSet$
alone and is unchanged: every $\mathbf q\in\acSet$ vanishes outside the
\tref{acceptance support} $\Chacc$ of \eqref{eq:chacc-def}, and for
$\mathbf q,\mathbf q'\in\acSet$ every feasible direction has
$q'(\hat c)-q(\hat c)=0$ at every $\hat c\notin\Chacc$. Part~(ii) holds for
$O^{J,\mathcal{D}}_{\alpha}$ on $\acSet\times\Feas^{J}(\varepsilon)$, for
every $\varepsilon\ge0$, with the cone clause of
Theorem~\ref{thm:branch-formula} replaced by
Proposition~\ref{prop:drop-bound}\,(i), which states
$0\le K^{\mathcal{D}}_{\hat c}\le P_{\hat c}$ for every $\omega$ block diagonal
along the joint sectors with positive semidefinite blocks; every point of
$\Feas^{J}(\varepsilon)$ is per-block positive semidefinite, the one
condition that set does not relax (Definition~\ref{def:relaxed-feasible-set}),
while $\chi$ preserves the joint-sector block structure and positivity
(Lemma~\ref{lem:choi-factor}\,(ii); in untrusted-source mode $\chi$ is the
identity). So for
$\hat c\notin\Chacc$, where $q(\hat c)=0$, the convention
$0\log_2\frac0v=0$ of \xref{foundations/def:kl}, stated for every
$v\ge0$ and the case $v=0$ included, makes the term
$\frac{\alpha}{\alpha-1}q(\hat c)\log_2\frac{q(\hat c)}{P^{J}_{\hat c}}$ and,
for $\hat c\in\Chkey$, the pair
$q(\hat c)\log_2\frac{q(\hat c)}{P^{J}_{\hat c}}
+\frac{q(\hat c)}{\alpha-1}\log_2\frac{q(\hat c)}{K^{J,\mathcal{D}}_{\hat c}}$
vanish at \emph{every} point of that set, whatever $P^{J}_{\hat c}$ and
$K^{J,\mathcal{D}}_{\hat c}$ are there and in particular when one of them
vanishes. Deleting those terms changes no value, so
$O^{J,\mathcal{D}}_{\alpha}$ is on that set a function of
$(\mathbf q|_{\Chacc},J)$ alone, (H-POS$^{\mathcal{D}}$) is stated on
$\Chacc$, and on $\Chkey\cap\Chacc$ for the $K$-part, exactly where it is
consumed, and a component of the acceptance-side gradient indexed outside
$\Chacc$ is paired against a direction component that is $0$ in every pairing
below, so its value is never read. No normalisation, coupling condition
or declared row is read in this step, which is why it holds on the
relaxed sets and not on the exact set alone.

\emph{Step 1 (the gradient exists and is a directional derivative in every
feasible direction).} Fix $t$, write $x\coloneqq x_t=(\mathbf q,J)$ and
$\nabla O\coloneqq(\lmoq_t,\lmoJ_t)$, let
$y=(\mathbf q',J')\in\acSet\times\Feas^{J}$ and $h\coloneqq y-x$, and put
$\Delta_p\coloneqq J'_p-J_{t,p}$, Hermitian as a difference of positive
semidefinite operators. By Step~0$'$ the objective along the segment is the
finite sum over $\hat c\in\Chacc$ of
$\frac{\alpha}{\alpha-1}q(\hat c)\log_2\frac{q(\hat c)}{P^{J}_{\hat c}}$ for
$\hat c\notin\Chkey$ and of the pair of Step~0$'$ for
$\hat c\in\Chkey\cap\Chacc$. In it $q(\hat c)$ and $P^{J}_{\hat c}$ are affine
and strictly positive at $x$ by (H-POS$^{\mathcal{D}}$), and
\begin{equation}\label{eq:KD-choi}
  K^{J,\mathcal{D}}_{\hat c}(J)=P^{J}_{\hat c}(J)
  +\sum_{p:\,(\hat c,p)\notin\mathcal{D}}
  f_{\hat c,p}\big(\chi_p(J_p)\big)
\end{equation}
by \eqref{eq:KD}, with $f_{\hat c,p}=T\circ\cG_{\hat c,p}
-\tr\circ\,\cG_{\hat c,p}$ and the second summand linear in $J_p$. Each kept
summand of \eqref{eq:KD-choi} falls into one of two cases. If $p$ is key-inert
for $\hat c$ then $f_{\hat c,p}\equiv0$ on the whole subcone
(Lemma~\ref{lem:joint-sectors}\,(iii)), hence identically along
the segment: the summand is constant, contributes $0$ to the derivative, and
no support hypothesis on it is consumed, which is why
(H-SUPP$^{\mathcal{D}}$) asks nothing of it. If $p\in\mathbb{J}_K(\hat c)$ and
$(\hat c,p)\notin\mathcal{D}$ then, by (H-SUPP$^{\mathcal{D}}$) with Step~0(c),
$\operatorname{supp}X_{\hat c,p}=W_{\hat c,p}$, and by Step~0(b)
$\operatorname{supp}\cG^{J}_{\hat c,p}(\Delta_p)\subseteq W_{\hat c,p}
=\operatorname{supp}X_{\hat c,p}$; so
Theorem~\ref{thm:gradient}\,(i) applies to the pair
$(X_{\hat c,p},\cG^{J}_{\hat c,p}(\Delta_p))$ and
$t'\mapsto T(X_{\hat c,p}+t'\cG^{J}_{\hat c,p}(\Delta_p))$ is differentiable at
$t'=0$, two-sidedly, the hypothesis being symmetric under
$\Delta_p\mapsto-\Delta_p$. And the $+\infty$ case of
Theorem~\ref{thm:gradient}\,(ii) is unreachable: that case sits
at directions strictly enlarging the support, Step~0(b) says no image of a
Hermitian direction leaves $W_{\hat c,p}$, and the directions this proof pairs
against are admissible for it because
$X_{\hat c,p}+t'\cG^{J}_{\hat c,p}(\Delta_p)
=\cG^{J}_{\hat c,p}\big((1-t')J_{t,p}+t'J'_p\big)\succeq0$ for $t'\in[0,1]$.
This is Lemma~\ref{lem:fo-face}\,(a), (b) read on the joint pair, with
Step~0(b) in place of \eqref{eq:fo-face}. A pair with
$(\hat c,p)\in\mathcal{D}$ is absent from \eqref{eq:KD-choi} altogether, so no
matrix power, no eigendecomposition and no support condition of a fully dropped
joint sector enters
(Corollary~\ref{cor:drop-program}\,(iii)). The remaining
dependence is through finitely many compositions of $\log_2$ with quantities
(H-POS$^{\mathcal{D}}$) keeps strictly positive, all differentiable there.
Hence $t'\mapsto O^{J,\mathcal{D}}_{\alpha}(x+t'h)$ is differentiable at
$t'=0$ with derivative $\langle\nabla O,h\rangle$, where, by the chain and sum
rules exactly as in the proof of
Corollary~\ref{cor:full-gradient}, whose maximal-support step is
Step~0(b) here, the acceptance-side component is \eqref{eq:grad-q} with
$K^{J,\mathcal{D}}_{\hat c}$ in place of $K^{J}_{\hat c}$, read as a vector in
$\mathbb{R}^{\Chacc}$, and the block of the feasible-side component at the
joint pair $p$ is \eqref{eq:grad-J} with the sums restricted to $\Chacc$,
respectively $\Chkey\cap\Chacc$, with $K^{J,\mathcal{D}}_{\hat c}$ in the
scalar prefactors, with $\Gamma^{(p)}_{\hat c}$ of
\eqref{eq:grad-pieces} read on the joint pair (the operators
$W_{\mathsf t}^{\dagger}M^{(\mathsf t)}_xW_{\mathsf t}\otimes F^{(s)}_v$), and
with the deficit piece $\Delta^{(p)}_{\hat c}$ carrying, in place of
$\mathbb 1[s\in\SecSet_K(\hat c)]$, the indicator
$\mathbb 1[p\in\mathbb{J}_K(\hat c)\text{ and }(\hat c,p)\notin\mathcal{D}]$.
That is the object (H-NOPERT$^{\mathcal{D}}$) names, and on a fully
dropped joint sector every deficit piece is absent, so the block is
$\chi_p^{\dagger}$ applied to \eqref{eq:drop-gradient}, which is what that
hypothesis says of such a sector.

\emph{Step 2 (convexity turns the derivative into an inequality).}
$O^{\mathcal{D}}_{\alpha}$ is jointly convex on $\Delta(\Chatset)\times\Feas$
with values in $[0,+\infty]$
(Corollary~\ref{cor:drop-program}\,(ii)), and
$O^{J,\mathcal{D}}_{\alpha}=O^{\mathcal{D}}_{\alpha}\circ(\idmap\times\chi)$
with $\chi$ \emph{linear}, so it is convex on the convex preimage
$\acSet\times\Feas^{J}$, a pullback along a linear map that needs no
injectivity, and $\Feas^{J}$ is convex and compact by
Theorem~\ref{thm:choi} in its general format, in untrusted-source mode
because it is $\Feas$ (Definition~\ref{def:feasible}). $x$ and $y$ lie in
that convex set and $O^{J,\mathcal{D}}_{\alpha}(x)$ is finite by
(H-POS$^{\mathcal{D}}$), so the three-chord inequality for a convex
extended-real function gives, exactly as in Step~2 of the proof of
Theorem~\ref{thm:firstorder}\,(ii),
\begin{equation}\label{eq:fo-subgrad-dropped}
  O^{J,\mathcal{D}}_{\alpha}(y)\;\ge\;O^{J,\mathcal{D}}_{\alpha}(x)
  +\langle\nabla O,\,y-x\rangle
  \qquad\text{for every }y\in\acSet\times\Feas^{J},
\end{equation}
the case $O^{J,\mathcal{D}}_{\alpha}(y)=+\infty$ included. Taking the infimum
over $y$: by Theorem~\ref{thm:choi} in its general format
$\chi(\Feas^{J})=\Feas$, in untrusted-source mode because $\chi$ is the
identity and $\Feas^{J}=\Feas$, and surjectivity is all a value needs, so
$\inf_{\acSet\times\Feas^{J}}O^{J,\mathcal{D}}_{\alpha}
=\inf_{\acSet\times\Feas}O^{\mathcal{D}}_{\alpha}
=\underline{h}^{\mathcal{D}}_{\alpha}$ by \eqref{eq:hlowD}; hence
$\underline{h}^{\mathcal{D}}_{\alpha}\ge O^{J,\mathcal{D}}_{\alpha}(x)
+\min_{y\in\acSet\times\Feas^{J}}\langle\nabla O,y-x\rangle$, the minimum
attained because $\acSet\times\Feas^{J}$ is compact and the map linear.

\emph{Steps 3 and 4.} The map $y\mapsto\langle\nabla O,y\rangle$ is a sum of
one term in each factor and $\acSet\times\Feas^{J}$ is a product set, so the
minimum separates over the product and \eqref{eq:lmo}, at the linear forms
$(\lmoq_t,\lmoJ_t)$ of this corollary, bounds each part below by the oracle
return; substituting gives
$\underline{h}^{\mathcal{D}}_{\alpha}\ge O^{J,\mathcal{D}}_{\alpha}(x_t)
+\langle\lmoq_t,v_t-\mathbf q_t\rangle+\langle\lmoJ_t,S_t-J_t\rangle$ for each
$t\in\Itset$. Where $b^{\mathrm A}_t$ or $b^{\mathrm F}_t$ stands in for a
pairing, \eqref{eq:lmo-bounds} bounds that part below by it instead, and
the same substitution gives the bracket with the replacement made, exactly
as in the proof of Theorem~\ref{thm:firstorder}\,(iv); where both stand in,
no oracle return enters. $\Itset$ is finite and non-empty, so the largest of those
right-hand sides, which is \eqref{eq:hlb} for
$O^{J,\mathcal{D}}_{\alpha}$, that is
$\underline{h}^{\mathrm{FO},\mathcal{D}}_{\alpha}$, is a lower bound too.
Finally $\underline{h}^{\mathcal{D}}_{\alpha}\ge0$ by
Proposition~\ref{prop:drop-bound}\,(iii), a statement about
the exact set, so $\max\{0,
\underline{h}^{\mathrm{FO},\mathcal{D}}_{\alpha}\}$ is the maximum of two lower
bounds on $\underline{h}^{\mathcal{D}}_{\alpha}$ and is one. This is the first
inequality of \eqref{eq:hlb-dropped}.

\emph{The implications, and the degenerate case.} (H-POS)
$\Rightarrow$ (H-POS$^{\mathcal{D}}$): $K^{J,\mathcal{D}}_{\hat c}\ge
K^{J}_{\hat c}$ by
Proposition~\ref{prop:drop-bound}\,(i), so
$K^{J}_{\hat c}>0$ gives $K^{J,\mathcal{D}}_{\hat c}>0$, and the other clauses
are identical. (H-SUPP) on joint blocks $\Rightarrow$
(H-SUPP$^{\mathcal{D}}$): with $J_s=\bigoplus_{\mathsf t}J_{t,(\mathsf t,s)}$,
Lemma~\ref{lem:joint-sectors}\,(i) gives
$X_{\hat c,s}=\bigoplus_{\mathsf t}X_{\hat c,(\mathsf t,s)}$ along mutually
orthogonal subspaces, so $\operatorname{supp}X_{\hat c,s}
=\bigoplus_{\mathsf t}\operatorname{supp}X_{\hat c,(\mathsf t,s)}$, while
$W_{\hat c,s}=\bigoplus_{\mathsf t}W_{\hat c,(\mathsf t,s)}$ by
\eqref{eq:Wcs-joint-split} with
$\operatorname{supp}X_{\hat c,(\mathsf t,s)}\subseteq
W_{\hat c,(\mathsf t,s)}$ by Step~0(c); two direct sums over the same
orthogonal summands that are equal with term-wise inclusions are equal term by
term, so maximal support of the sector gives maximal support of each of its
joint pairs. At $\mathcal{D}=\varnothing$ and $|\mathbb{T}|=1$:
$K^{\mathcal{D}}_{\hat c}=K_{\hat c}$ and
$O^{\mathcal{D}}_{\alpha}=O_{\alpha}$ by
Definition~\ref{def:dropped},
$\mathbb{J}_K(\hat c)=\{t_0\}\times\SecSet_K(\hat c)$ and
$W_{\hat c,(t_0,s)}=W_{\hat c,s}$ by
Definition~\ref{def:joint-sectors} and
\eqref{eq:Wcs-joint-split}, so the three hypotheses are (H-POS), (H-SUPP),
(H-NOPERT) and \eqref{eq:hlb-dropped} is \eqref{eq:hlb} together with
$\hlow\le\hqkd(\CList)$: Theorem~\ref{thm:firstorder}\,(ii).

\emph{The relaxed reading.} The proof is the one above with a single change,
made in Step~2: $x$ ranges over $\acSet\times\Feas^{J}(\varepsilon)$ while the
comparison point $y$ ranges over the \emph{exact} $\acSet\times\Feas^{J}$.
Steps~0, 0$'$ and~1 are unchanged and read only per-block positive
semidefiniteness of $x$: Step~0(b), (c) ask that of $J_{t,p}$, Step~0$'$
asks it through
Proposition~\ref{prop:drop-bound}\,(i), and $\Delta_p$ is
Hermitian as a difference of a positive semidefinite $J'_p$ from $\Feas^{J}$
and a positive semidefinite $J_{t,p}$ from $\Feas^{J}(\varepsilon)$, per-block
positive semidefiniteness being the one condition that set does not relax
(Definition~\ref{def:relaxed-feasible-set}). In Step~2 both endpoints lie in
$\acSet\times\Feas^{J}(\varepsilon)$, which is convex
(Proposition~\ref{prop:relaxed-set-properties}\,(b); $\acSet$ by
Definition~\ref{def:admissible-acceptance-set}), and
$O^{J,\mathcal{D}}_{\alpha}$ is convex on it. The transfer is the one
Proposition~\ref{prop:relaxed-set-properties}\,(c) performs for
Theorem~\ref{thm:convexity}, and it goes through for the dropped
objective for the same reason: the convexity argument of
Corollary~\ref{cor:drop-program}\,(ii) reads of its domain only
that it is a convex set on which $P_{\hat c}$ is affine and
$K^{\mathcal{D}}_{\hat c}$ is non-negative, concave and continuous;
Corollary~\ref{cor:drop-program}\,(i) states that concavity and
continuity on the whole \emph{cone} of operators block diagonal along the joint
sectors, while
Proposition~\ref{prop:drop-bound}\,(i) states
$0\le K^{\mathcal{D}}_{\hat c}\le P_{\hat c}$ there; and
$\chi(\Feas^{J}(\varepsilon))$ is a convex subset of that cone, so
$O^{\mathcal{D}}_{\alpha}$ is convex on
$\acSet\times\chi(\Feas^{J}(\varepsilon))$ and pulls back along the linear
$\chi$. So the three-chord inequality gives
\eqref{eq:fo-subgrad-dropped} for every $y\in\acSet\times\Feas^{J}$, and the
infimum is taken over \emph{that} set, not over the relaxed one; Steps~3
and~4 are then as above, and the clamp reads
Proposition~\ref{prop:drop-bound}\,(iii) at $\varepsilon=0$,
where it holds. What this uses of the relaxed set, in full: that
$\Feas^{J}\subseteq\Feas^{J}(\varepsilon)$
(Proposition~\ref{prop:relaxed-set-properties}\,(a)); that
$\acSet\times\Feas^{J}(\varepsilon)$ is convex (\emph{ibid.}\,(b)); that every
member is per-block positive semidefinite; and the convexity transfer just
made, whose own derivation consumes the compactness of (b) through the
continuity of $K^{\mathcal{D}}_{\hat c}$, so compactness is used one level down
and never in the minimisation. Of the \emph{exact} set it uses compactness,
non-emptiness, and that the infimum of $O^{J,\mathcal{D}}_{\alpha}$ over it is
$\underline{h}^{\mathcal{D}}_{\alpha}$. No statement is made about the
infimum of $O^{J,\mathcal{D}}_{\alpha}$ over the relaxed set, and none is
used.

\emph{An empty feasible set.} No claim is used for the two values being
$+\infty$. By
\eqref{eq:hlowD} and \eqref{eq:hlow-def}, $\underline{h}^{\mathcal{D}}_{\alpha}$
and $\hlow$ are infima over $\mathbf q\in\acSet$ and $\omega\in\Feas$; with
$\Feas=\varnothing$ each inner infimum is over the empty set and is $+\infty$
for every $\mathbf q$, and so is the outer one, whether $\acSet$ is empty or
not. Every real number is below $+\infty$. For the consequence: where the
tuple satisfies condition~(vii) as well, it satisfies
Definition~\ref{def:admissible-declaration} in full, and with it the two
conditions Definition~\ref{def:declaration} places on its raw-key map, which
are conditions (v) and (viii); it is then a protocol declaration satisfying
the hypothesis of the statement above; where it does not, $\Feas=\varnothing$, both values
are $+\infty$, and a real number is below them.
\end{proof}
